\chapter{Data Structures --- Linked Lists and Heaps}
\label{ch:linked_lists_and_heaps}

This chapter covers three major linked-list structures---singly linked lists, doubly linked lists, and circular linked lists---followed by heaps, with particular emphasis on step-by-step max-heap construction, insertion, deletion, and complexity analysis.

The explanations and implementations are aimed at Bangladesh government IT recruitment examinations, including Assistant Programmer, AD (ICT), Assistant Network Engineer, Assistant Maintenance Engineer, and Assistant Database Administrator positions.

\begin{center}
\vspace{0.5em}
{\color{primary}\rule{0.45\textwidth}{1.2pt}}\\[4pt]
{\Large\bfseries\color{primary} Section I: Linked Lists}\\[2pt]
{\color{primary}\rule{0.45\textwidth}{1.2pt}}
\vspace{0.5em}
\end{center}

\section{What is a Linked List?}

A \textbf{linked list} is a linear data structure consisting of a sequence of nodes. Each node stores data and one or more links (pointers) to other nodes.

Unlike an array, whose elements are conventionally stored in contiguous memory locations, the nodes of a linked list can occupy different memory locations across the heap. Pointers connect them into a coherent logical sequence.

A node in a singly linked list contains two fields:
\begin{itemize}
    \item \textbf{Data:} The value stored in the node.
    \item \textbf{Next pointer:} The memory address of the next node in the sequence.
\end{itemize}

For example, consider a list containing $10$, $20$, and $30$. The first node stores $10$ and points to the node containing $20$; the second points to the node containing $30$; the final node's next pointer is \texttt{NULL} (or \texttt{nullptr}) in a conventional, non-circular singly linked list:
\[
\text{Head} \longrightarrow [10 \mid \bullet] \longrightarrow [20 \mid \bullet] \longrightarrow [30 \mid \text{nullptr}]
\]

\subsection{Important Terminology}

\begin{table}[htbp]
\centering
\small
\begin{tabularx}{\textwidth}{l X}
\toprule
\textbf{Term} & \textbf{Meaning} \\
\midrule
Node                      & An individual unit containing data and pointer/link fields \\
Head                      & A pointer or reference addressing the first node in the list \\
Tail                      & A pointer or reference addressing the last node, if maintained \\
Link                      & A pointer or reference connecting one node to another \\
\texttt{NULL} / \texttt{nullptr} & A sentinel value indicating the absence of a successor node \\
Traversal                 & Visiting each node sequentially from head to end \\
Insertion                 & Allocating and linking a new node into the list \\
Deletion                  & Unlinking and deallocating an existing node \\
Dynamic Memory Allocation & Allocating node memory at runtime as needed \\
\bottomrule
\end{tabularx}
\caption{Fundamental linked list terminology.}
\end{table}

In C++, nodes are commonly created dynamically using \texttt{new} and deallocated using \texttt{delete}. In C, \texttt{malloc()} and \texttt{free()} are standard.

\section{Why Use Linked Lists Instead of Arrays?}

An array is usually preferable when fast indexed access ($O(1)$ random access) is important. A linked list is advantageous when the application performs frequent insertions or deletions at positions where the target node or link is already known.

Consider an array containing:
\[
[10, \; 20, \; 30, \; 40]
\]
Inserting $25$ between $20$ and $30$ requires shifting all subsequent elements ($30$ and $40$) to the right, incurring $O(n)$ time.

In a linked list, if a pointer to the node containing $20$ is already available, a new node containing $25$ can be inserted in $O(1)$ time by adjusting two pointers. Existing nodes remain stationary in memory. However, locating the node containing $20$ from the head may itself require an $O(n)$ linear traversal.

\subsection{Array versus Linked List Comparison}

\begin{table}[htbp]
\centering
\footnotesize
\begin{tabularx}{\textwidth}{l X X}
\toprule
\textbf{Basis} & \textbf{Array} & \textbf{Linked List} \\
\midrule
Memory Layout                   & Contiguous memory block                           & Scattered nodes linked via pointers \\
Size Flexibility                & Fixed (static) or reallocation cost (dynamic)      & Fully dynamic; grows/shrinks per node \\
Random Access                   & $O(1)$ via index arithmetic                       & $O(n)$ sequential link traversal \\
Access to First Element         & $O(1)$                                            & $O(1)$ via \texttt{head} pointer \\
Insert at Beginning             & $O(n)$ due to element shifting                    & $O(1)$ pointer reassignment \\
Delete at Beginning             & $O(n)$ due to element shifting                    & $O(1)$ pointer reassignment \\
Insert After Known Node         & $O(n)$ due to element shifting                    & $O(1)$ pointer updates \\
Search in Unsorted Data         & $O(n)$                                            & $O(n)$ \\
Extra Memory Overhead           & Minimal structural overhead                       & Extra memory per node for pointers \\
Cache Locality                  & Excellent (spatial locality)                      & Generally poor (cache misses common) \\
Binary Search Capability        & $O(\log n)$ on sorted array                       & Infeasible for $O(\log n)$ (middle node is not $O(1)$) \\
\bottomrule
\end{tabularx}
\caption{Architectural comparison between Arrays and Linked Lists.}
\end{table}

\begin{quote}
\textbf{Important Exam Distinction:} Inserting a node into a linked list takes $O(1)$ time \emph{if the predecessor pointer is already held}. Finding that position from scratch still takes $O(n)$ traversal time.
\end{quote}

\section{Singly Linked List}

\subsection{Definition and Node Structure}

A \textbf{singly linked list (SLL)} is a linked list in which each node contains data and exactly one pointer to the next node. Each node can directly identify its successor, but cannot directly identify its predecessor.

The first node is accessed through \texttt{head}. The last node is identified by \texttt{next == nullptr}.

\begin{lstlisting}[style=cppstyle]
struct Node {
    int data;
    Node* next;

    Node(int value) : data(value), next(nullptr) {}
};
\end{lstlisting}

\subsubsection{Main Characteristics}
\begin{enumerate}
    \item Each node has exactly one link field.
    \item Traversal proceeds strictly forward.
    \item Nodes are inserted or removed without shifting remaining elements in memory.
    \item Reverse traversal is not directly supported without auxiliary structures (e.g., stack) or pointer reversal.
    \item Accessing the $k$-th element requires $O(k)$ pointer dereferences from \texttt{head}.
\end{enumerate}

\subsection{Traversing a Singly Linked List}

Traversal visits every node in order from \texttt{head} until \texttt{nullptr} is reached.

\begin{lstlisting}[style=cppstyle]
void display(Node* head) {
    Node* current = head;
    while (current != nullptr) {
        cout << current->data << " ";
        current = current->next;
    }
    cout << '\n';
}
\end{lstlisting}

\textbf{Complexity:} Time is $O(n)$ as all $n$ nodes are visited once; auxiliary space is $O(1)$.

\subsection{Insertion Operations}

\subsubsection{A. Insert at the Beginning (Front)}
To insert value $x$ before the current head:
\begin{enumerate}
    \item Allocate a new node containing $x$.
    \item Point \texttt{newNode->next} to the current \texttt{head}.
    \item Reassign \texttt{head} to \texttt{newNode}.
\end{enumerate}

\begin{lstlisting}[style=cppstyle]
void insertFront(Node*& head, int value) {
    Node* newNode = new Node(value);
    newNode->next = head;
    head = newNode;
}
\end{lstlisting}

\textbf{Time Complexity:} $O(1)$. If the list is initially empty (\texttt{head == nullptr}), \texttt{newNode->next} becomes \texttt{nullptr}, and \texttt{newNode} correctly becomes the head.

\subsubsection{B. Insert at the End (Tail)}
To append a node at the tail:
\begin{enumerate}
    \item Allocate the new node with \texttt{next = nullptr}.
    \item If \texttt{head == nullptr}, set \texttt{head = newNode}.
    \item Otherwise, traverse from \texttt{head} to the node where \texttt{current->next == nullptr}.
    \item Assign \texttt{current->next = newNode}.
\end{enumerate}

\begin{lstlisting}[style=cppstyle]
void insertEnd(Node*& head, int value) {
    Node* newNode = new Node(value);
    if (head == nullptr) {
        head = newNode;
        return;
    }
    Node* current = head;
    while (current->next != nullptr) {
        current = current->next;
    }
    current->next = newNode;
}
\end{lstlisting}

\textbf{Time Complexity:} $O(n)$ when only \texttt{head} is tracked. If an explicit \texttt{tail} pointer is maintained, insertion at the end is $O(1)$.

\subsubsection{C. Insert After a Specified Node}
Suppose \texttt{current} points to the node after which the new node is to be inserted:
\begin{enumerate}
    \item Allocate \texttt{newNode}.
    \item Set \texttt{newNode->next = current->next}.
    \item Set \texttt{current->next = newNode}.
\end{enumerate}

\begin{lstlisting}[style=cppstyle]
void insertAfter(Node* current, int value) {
    if (current == nullptr) return;
    Node* newNode = new Node(value);
    newNode->next = current->next;
    current->next = newNode;
}
\end{lstlisting}

\begin{quote}
\textbf{Common Exam Trap:} Always set \texttt{newNode->next} \emph{before} overwriting \texttt{current->next}. Overwriting \texttt{current->next} first breaks the link to the remainder of the list, causing memory leaks and orphaned nodes.
\end{quote}

\subsection{Deletion Operations}

\subsubsection{A. Delete the First Node (Front)}
\begin{lstlisting}[style=cppstyle]
void deleteFront(Node*& head) {
    if (head == nullptr) return;
    Node* temp = head;
    head = head->next;
    delete temp;
}
\end{lstlisting}
\textbf{Time Complexity:} $O(1)$.

\subsubsection{B. Delete a Node Containing a Target Value}
Deleting a value requires maintaining two pointers: \texttt{current} (the node being examined) and \texttt{previous} (the node preceding \texttt{current}).

\begin{lstlisting}[style=cppstyle]
void deleteValue(Node*& head, int value) {
    Node* current = head;
    Node* previous = nullptr;

    while (current != nullptr && current->data != value) {
        previous = current;
        current = current->next;
    }

    if (current == nullptr) return; // Value not found

    if (previous == nullptr) {
        head = current->next; // Target is head
    } else {
        previous->next = current->next; // Bypass current
    }
    delete current;
}
\end{lstlisting}

\textbf{Time Complexity:} $O(n)$ worst-case to locate the element.

\subsection{Advantages and Disadvantages of Singly Linked Lists}

\textbf{Advantages:}
\begin{itemize}
    \item Dynamic sizing without pre-allocation overhead.
    \item $O(1)$ insertion and deletion at the front.
    \item No need to shift memory elements upon modification.
    \item Simpler pointer management and lower memory overhead than doubly linked lists.
\end{itemize}

\textbf{Disadvantages:}
\begin{itemize}
    \item Unidirectional traversal only.
    \item No $O(1)$ random access.
    \item Deleting an arbitrary known node requires locating its predecessor ($O(n)$ search).
    \item Extra pointer memory per node ($4$ or $8$ bytes).
    \item Unfriendly to CPU cache lines due to non-contiguous heap allocations.
\end{itemize}

\subsection{Applications}
\begin{itemize}
    \item Implementing dynamic stacks and queues.
    \item Representing adjacency lists in graph algorithms.
    \item Separate chaining collision resolution in hash tables.
    \item Polynomial manipulation and arbitrary-precision arithmetic.
\end{itemize}

\section{Doubly Linked List}

\subsection{Definition and Node Structure}

A \textbf{doubly linked list (DLL)} is a linked list in which each node contains:
\begin{enumerate}
    \item Data payload.
    \item A pointer to the previous node (\texttt{prev}).
    \item A pointer to the next node (\texttt{next}).
\end{enumerate}

The first node has \texttt{prev == nullptr}, and the last node has \texttt{next == nullptr}.

\begin{lstlisting}[style=cppstyle]
struct DNode {
    int data;
    DNode* prev;
    DNode* next;

    DNode(int value) : data(value), prev(nullptr), next(nullptr) {}
};
\end{lstlisting}

\subsection{Singly versus Doubly Linked Lists}

\textbf{Advantages of DLL over SLL:}
\begin{itemize}
    \item Supports bidirectional traversal (forward and backward).
    \item Deletion of a known node is $O(1)$ because the predecessor is accessible directly via \texttt{node->prev}.
    \item Inserting before a known node is $O(1)$.
\end{itemize}

\textbf{Disadvantages of DLL:}
\begin{itemize}
    \item Requires two pointer fields per node, increasing memory overhead.
    \item Every insertion and deletion requires updating up to four pointers, increasing code complexity and bug potential.
\end{itemize}

\subsection{Insertion Operations}

\subsubsection{A. Insert at the Beginning}
\begin{lstlisting}[style=cppstyle]
void insertFront(DNode*& head, int value) {
    DNode* newNode = new DNode(value);
    newNode->next = head;
    if (head != nullptr) {
        head->prev = newNode;
    }
    head = newNode;
}
\end{lstlisting}
\textbf{Time Complexity:} $O(1)$.

\subsubsection{B. Insert After a Known Node}
\begin{lstlisting}[style=cppstyle]
void insertAfter(DNode* current, int value) {
    if (current == nullptr) return;
    DNode* newNode = new DNode(value);

    newNode->next = current->next;
    newNode->prev = current;

    if (current->next != nullptr) {
        current->next->prev = newNode;
    }
    current->next = newNode;
}
\end{lstlisting}
\textbf{Time Complexity:} $O(1)$ when \texttt{current} is given.

\subsection{Deletion of a Known Node}

A major advantage of a DLL is unlinking a known node in $O(1)$ without needing to traverse from the head:

\begin{lstlisting}[style=cppstyle]
void deleteNode(DNode*& head, DNode* node) {
    if (node == nullptr) return;

    if (node->prev != nullptr) {
        node->prev->next = node->next;
    } else {
        head = node->next; // node was head
    }

    if (node->next != nullptr) {
        node->next->prev = node->prev;
    }
    delete node;
}
\end{lstlisting}
\textbf{Time Complexity:} $O(1)$.

\subsection{Applications of Doubly Linked Lists}
\begin{itemize}
    \item Web browser navigation history (Back and Forward buttons).
    \item Undo and Redo functionality in editors.
    \item Least Recently Used (LRU) Cache implementation (combined with a hash map).
    \item Media player playlists with Next and Previous track navigation.
    \item Operating system thread and process schedulers.
\end{itemize}

\section{Circular Linked List}

\subsection{Definition and Principal Variants}

A \textbf{circular linked list} is a linked list in which the last node links back to the first node instead of holding \texttt{nullptr}. There are two variants:
\begin{enumerate}
    \item \textbf{Circular Singly Linked List (CSLL):} The last node's \texttt{next} points to \texttt{head}.
    \item \textbf{Circular Doubly Linked List (CDLL):} The last node's \texttt{next} points to \texttt{head}, and \texttt{head->prev} points to the last node.
\end{enumerate}

\subsection{Traversing a Circular Singly Linked List}

Because a circular list contains no \texttt{nullptr}, a standard \texttt{while (current != nullptr)} loop will run indefinitely. Instead, use a \texttt{do-while} loop:

\begin{lstlisting}[style=cppstyle]
void displayCircular(Node* head) {
    if (head == nullptr) return;
    Node* current = head;
    do {
        cout << current->data << " ";
        current = current->next;
    } while (current != head);
    cout << '\n';
}
\end{lstlisting}

\subsection{Insertion Using a Tail Pointer}

Maintaining a pointer to the \textbf{last node (\texttt{tail})} provides $O(1)$ insertion at both the head and the tail because:
\[
\text{Head} = \texttt{tail->next}
\]

\begin{lstlisting}[style=cppstyle]
class CircularLinkedList {
private:
    struct Node {
        int data;
        Node* next;
        Node(int val) : data(val), next(nullptr) {}
    };
    Node* tail = nullptr;

public:
    // Insert at front: O(1)
    void insertFront(int value) {
        Node* newNode = new Node(value);
        if (tail == nullptr) {
            newNode->next = newNode;
            tail = newNode;
            return;
        }
        newNode->next = tail->next;
        tail->next = newNode;
    }

    // Insert at end: O(1)
    void insertEnd(int value) {
        Node* newNode = new Node(value);
        if (tail == nullptr) {
            newNode->next = newNode;
            tail = newNode;
            return;
        }
        newNode->next = tail->next;
        tail->next = newNode;
        tail = newNode; // Only difference: tail moves to newNode
    }

    // Delete front: O(1)
    void deleteFront() {
        if (tail == nullptr) return;
        Node* head = tail->next;
        if (head == tail) {
            delete head;
            tail = nullptr;
            return;
        }
        tail->next = head->next;
        delete head;
    }

    void display() const {
        if (tail == nullptr) return;
        Node* current = tail->next;
        do {
            cout << current->data << " ";
            current = current->next;
        } while (current != tail->next);
        cout << '\n';
    }

    ~CircularLinkedList() {
        while (tail != nullptr) deleteFront();
    }
};
\end{lstlisting}

\subsection{Applications of Circular Linked Lists}
\begin{itemize}
    \item \textbf{Round-Robin CPU Scheduling:} Processes are allocated fixed time-slices in a cyclic rotation.
    \item \textbf{Continuous Media Streaming:} Repeating audio playlists or slideshow loops.
    \item \textbf{Multiplayer Turn-Based Games:} Cycling through players continuously.
\end{itemize}

\section{Comparison of Linked-List Types}

\begin{table}[htbp]
\centering
\small
\begin{tabularx}{\textwidth}{l X X X}
\toprule
\textbf{Feature} & \textbf{Singly Linked List} & \textbf{Doubly Linked List} & \textbf{Circular Singly Linked List} \\
\midrule
Links per Node       & One (\texttt{next})         & Two (\texttt{prev}, \texttt{next}) & One (\texttt{next}) \\
Traversal Direction  & Forward only                & Bidirectional               & Forward (cyclic) \\
End of List          & \texttt{next == nullptr}    & \texttt{next == nullptr}    & \texttt{next == head} \\
Insert at Head       & $O(1)$                      & $O(1)$                      & $O(1)$ (with tail pointer) \\
Insert at End        & $O(n)$ ($O(1)$ with tail)   & $O(n)$ ($O(1)$ with tail)   & $O(1)$ (with tail pointer) \\
Delete First Node    & $O(1)$                      & $O(1)$                      & $O(1)$ (with tail pointer) \\
Delete Known Node    & Requires predecessor ($O(n)$) & $O(1)$ unlinking           & Requires predecessor ($O(n)$) \\
Stopping Condition   & Reach \texttt{nullptr}      & Reach \texttt{nullptr}      & Return to starting node \\
Pointer Overhead     & 1 pointer per node          & 2 pointers per node         & 1 pointer per node \\
\bottomrule
\end{tabularx}
\caption{Comprehensive comparison of linked-list variants.}
\end{table}

\begin{center}
\vspace{1em}
{\color{primary}\rule{0.45\textwidth}{1.2pt}}\\[4pt]
{\Large\bfseries\color{primary} Section II: Heap Data Structures}\\[2pt]
{\color{primary}\rule{0.45\textwidth}{1.2pt}}
\vspace{0.5em}
\end{center}

\section{What is a Heap?}

A \textbf{heap} is a specialized tree-based data structure that satisfies two fundamental properties:
\begin{enumerate}
    \item \textbf{Shape Property:} The tree is a \textbf{complete binary tree}.
    \item \textbf{Heap Property:} The keys satisfy either the \textbf{max-heap} or \textbf{min-heap} ordering rule.
\end{enumerate}

\subsection{What is a Complete Binary Tree?}

A binary tree is \textbf{complete} if:
\begin{itemize}
    \item Every level, except possibly the last, is completely filled with nodes.
    \item All nodes at the last level reside as far left as possible, without gaps.
\end{itemize}

This property enables a binary heap to be represented compactly inside a contiguous 1D array without storing explicit parent or child pointers.

\begin{table}[htbp]
\centering
\small
\begin{tabularx}{\textwidth}{l X}
\toprule
\textbf{Binary Tree Type} & \textbf{Formal Definition} \\
\midrule
Complete Binary Tree & All levels full except possibly the last, which is filled strictly left-to-right \\
Full Binary Tree     & Every internal node has exactly two children (no node has 1 child) \\
Perfect Binary Tree  & All internal nodes have two children and all leaves reside at identical depth \\
\bottomrule
\end{tabularx}
\caption{Distinguishing Complete, Full, and Perfect binary trees.}
\end{table}

\subsection{Max-Heap vs Min-Heap}

\begin{itemize}
    \item \textbf{Max-Heap:} For every node $i$ other than the root:
    \[
    A[\text{parent}(i)] \geq A[i]
    \]
    The maximum element resides perpetually at the root ($A[0]$).
    \item \textbf{Min-Heap:} For every node $i$ other than the root:
    \[
    A[\text{parent}(i)] \leq A[i]
    \]
    The minimum element resides perpetually at the root ($A[0]$).
\end{itemize}

\subsection{Is a Heap the Same as a Binary Search Tree (BST)?}

No. This is a common viva question:
\begin{table}[htbp]
\centering
\small
\begin{tabularx}{\textwidth}{l X X}
\toprule
\textbf{Basis} & \textbf{Binary Heap} & \textbf{Binary Search Tree (BST)} \\
\midrule
Ordering Rule     & Vertical ordering: Parent $\geq$ children (no left vs right rule) & Horizontal ordering: Left $<$ Root $<$ Right \\
Root Node         & Always the global Maximum (max-heap) or Minimum (min-heap)        & Arbitrary key depending on insertion \\
Access to Extremum & $O(1)$                                                            & $O(h)$ ($O(\log n)$ balanced, $O(n)$ skewed) \\
Search for Arbitrary Key & $O(n)$ linear scan (no directional guarantee)                & $O(\log n)$ average / balanced \\
Underlying Storage & Implicit 1D Array (zero pointer overhead)                         & Explicit linked nodes with left/right pointers \\
\bottomrule
\end{tabularx}
\caption{Key distinctions between Binary Heap and Binary Search Tree.}
\end{table}

\section{Representing a Binary Heap Using an Array}

Because a heap is a complete binary tree, parent-child relationships map directly to arithmetic indices.

\subsection{Zero-Based Indexing Formulas}
For a node located at index $i$ in a zero-indexed array ($i = 0, 1, \dots, n-1$):
\begin{align*}
\text{Parent}(i) &= \left\lfloor \frac{i - 1}{2} \right\rfloor \quad (\text{for } i > 0) \\
\text{Left Child}(i) &= 2i + 1 \\
\text{Right Child}(i) &= 2i + 2
\end{align*}

\subsection{One-Based Indexing Formulas}
For a one-indexed array ($i = 1, 2, \dots, n$):
\begin{align*}
\text{Parent}(i) &= \left\lfloor \frac{i}{2} \right\rfloor \\
\text{Left Child}(i) &= 2i \\
\text{Right Child}(i) &= 2i + 1
\end{align*}

\subsubsection{Example Mapping}
Consider the array $[50, 30, 40, 10, 20]$ of size $n=5$:
\begin{table}[htbp]
\centering
\small
\begin{tabularx}{0.85\textwidth}{ccccc}
\toprule
\textbf{Index $i$} & \textbf{Value $A[i]$} & \textbf{Parent Index} & \textbf{Left Child ($2i+1$)} & \textbf{Right Child ($2i+2$)} \\
\midrule
0 & 50 & None & 1 (value 30) & 2 (value 40) \\
1 & 30 & 0    & 3 (value 10) & 4 (value 20) \\
2 & 40 & 0    & 5 (out of range) & 6 (out of range) \\
3 & 10 & 1    & --- & --- \\
4 & 20 & 1    & --- & --- \\
\bottomrule
\end{tabularx}
\caption{Zero-based array index mapping of a binary heap.}
\end{table}

\section{Max-Heap Construction: Step by Step}

Given an arbitrary array, two strategies exist to convert it into a valid max-heap:
\begin{enumerate}
    \item \textbf{Repeated Insertion:} Insert elements one by one, sifting up ($O(n \log n)$ time).
    \item \textbf{Bottom-Up Heap Construction (Floyd's Algorithm):} Place elements into array, then call \texttt{maxHeapify} downward starting from the last non-leaf node up to the root ($O(n)$ time).
\end{enumerate}

\subsection{The Heapify (Sift-Down) Operation}

\texttt{maxHeapify(A, i, n)} assumes that the subtrees rooted at $\text{left}(i)$ and $\text{right}(i)$ are valid max-heaps, but $A[i]$ may violate the heap property. It floats $A[i]$ downward by repeatedly swapping with its largest child until the property is restored.

\subsection{Step-by-Step Bottom-Up Construction Trace}

Given the input array:
\[
A = [4, \; 10, \; 3, \; 5, \; 1], \quad n = 5
\]
The last non-leaf node index in a zero-indexed array is:
\[
\text{Last non-leaf} = \left\lfloor \frac{n}{2} \right\rfloor - 1 = \left\lfloor \frac{5}{2} \right\rfloor - 1 = 1
\]
Leaves reside at indices $2, 3, 4$ and are trivial heaps. We only process indices $1$ down to $0$.

\begin{enumerate}
    \item \textbf{Step 1: Heapify Node at Index 1 ($A[1] = 10$)}
    \begin{itemize}
        \item Left child at index $3$: $A[3] = 5$.
        \item Right child at index $4$: $A[4] = 1$.
        \item Since $10 \geq 5$ and $10 \geq 1$, node $1$ already satisfies max-heap ordering. No swap.
        \item Array remains: $[4, 10, 3, 5, 1]$.
    \end{itemize}

    \item \textbf{Step 2: Heapify Node at Index 0 ($A[0] = 4$)}
    \begin{itemize}
        \item Left child at index $1$: $A[1] = 10$.
        \item Right child at index $2$: $A[2] = 3$.
        \item Largest among $\{4, 10, 3\}$ is $10$ at index $1$.
        \item Swap $A[0]$ and $A[1]$.
        \item Array becomes: $[10, 4, 3, 5, 1]$.
    \end{itemize}

    \item \textbf{Step 3: Continue Sifting Down Displaced Node ($4$ at Index 1)}
    \begin{itemize}
        \item Children of index $1$ are $A[3] = 5$ and $A[4] = 1$.
        \item Largest among $\{4, 5, 1\}$ is $5$ at index $3$.
        \item Swap $A[1]$ and $A[3]$.
        \item Array becomes: $[10, 5, 3, 4, 1]$.
        \item Node $4$ is now at index $3$ (a leaf), so sift-down terminates.
    \end{itemize}
\end{enumerate}

\textbf{Final Valid Max-Heap:} $[10, 5, 3, 4, 1]$.

\begin{table}[htbp]
\centering
\small
\begin{tabularx}{0.85\textwidth}{cc X c}
\toprule
\textbf{Step} & \textbf{Index Processed} & \textbf{Action Taken} & \textbf{Array State} \\
\midrule
Initial & --- & Original sequence & $[4, 10, 3, 5, 1]$ \\
1       & 1   & Checked children $\{5, 1\}$; no swap & $[4, 10, 3, 5, 1]$ \\
2       & 0   & Swapped $4 \leftrightarrow 10$ with child & $[10, 4, 3, 5, 1]$ \\
3       & 1   & Swapped $4 \leftrightarrow 5$ to restore leaf & $[10, 5, 3, 4, 1]$ \\
\bottomrule
\end{tabularx}
\caption{Trace of bottom-up max-heap construction.}
\end{table}

\subsection{Why Does Bottom-Up Heap Construction Take $O(n)$?}

A single \texttt{heapify} takes $O(\log n)$ worst-case time. However, bottom-up construction does not run for full height on all nodes:
\begin{itemize}
    \item Half the nodes are leaves ($h=0$), requiring $0$ operations.
    \item A quarter of the nodes have height $1$, moving down at most $1$ step.
    \item Only $1$ node (the root) can travel the full tree height $\lfloor \log n \rfloor$.
\end{itemize}

Summing work over all tree heights:
\[
T(n) = \sum_{h=0}^{\lfloor \log n \rfloor} \left\lceil \frac{n}{2^{h+1}} \right\rceil O(h) = O\left( n \sum_{h=0}^{\infty} \frac{h}{2^h} \right)
\]
Since the infinite series $\sum_{h=0}^{\infty} \frac{h}{2^h} = 2$, we obtain:
\[
T(n) = O(2n) = O(n)
\]
This proves that bottom-up construction runs in \textbf{linear time $O(n)$}.

\section{Max-Heap Insertion}

Insertion maintains both the complete binary tree shape and the max-heap ordering via \textbf{sift-up} (percolate-up).

\subsection{Algorithm}
\begin{enumerate}
    \item Append the new value $x$ at the end of the array ($A[n]$), expanding the heap size.
    \item Compare $x$ with its parent: $p = \lfloor (i - 1)/2 \rfloor$.
    \item If $x > A[p]$, swap $A[i]$ and $A[p]$, and set $i = p$.
    \item Repeat until $x \leq A[p]$ or $x$ reaches the root ($i = 0$).
\end{enumerate}

\subsection{Step-by-Step Insertion Example}
Insert $15$ into existing max-heap $[10, 5, 3, 4, 1]$:
\begin{enumerate}
    \item Append $15$ at index $5$: $[10, 5, 3, 4, 1, \mathbf{15}]$.
    \item Parent of index $5$ is $\lfloor(5-1)/2\rfloor = 2$ ($A[2] = 3$). Since $15 > 3$, swap:
    \[
    [10, 5, \mathbf{15}, 4, 1, \mathbf{3}]
    \]
    \item New parent of index $2$ is $\lfloor(2-1)/2\rfloor = 0$ ($A[0] = 10$). Since $15 > 10$, swap:
    \[
    [\mathbf{15}, 5, \mathbf{10}, 4, 1, 3]
    \]
    \item Index $0$ is the root. Insertion terminates.
\end{enumerate}
\textbf{Complexity:} Worst-case time is $O(\log n)$ (tree height). Auxiliary space is $O(1)$.

\section{Max-Heap Deletion (Extract Maximum)}

Deleting from a max-heap refers to extracting the root node containing the maximum key.

\subsection{Algorithm}
\begin{enumerate}
    \item Save root $A[0]$ for return.
    \item Copy the last element in the array to the root ($A[0] = A[n-1]$).
    \item Decrement heap size by $1$ (removing the last position).
    \item Call \texttt{maxHeapify(A, 0, n)} to sift the new root downward until the max-heap property is restored.
\end{enumerate}

\subsection{Step-by-Step Deletion Example}
Extract max from $[15, 5, 10, 4, 1, 3]$:
\begin{enumerate}
    \item Save maximum: $\text{max} = 15$.
    \item Replace root with last element $3$: $[3, 5, 10, 4, 1]$.
    \item Children of root ($3$) are $A[1] = 5$ and $A[2] = 10$. Largest child is $10$ at index $2$.
    \item Swap $3$ and $10$:
    \[
    [10, 5, \mathbf{3}, 4, 1]
    \]
    \item Node $3$ at index $2$ has no children ($2(2)+1 = 5 \geq 5$). Sift-down terminates.
\end{enumerate}
\textbf{Complexity:} $O(\log n)$ worst-case time.

\subsection{Deleting an Arbitrary Key}
To delete an arbitrary key $k$ (not necessarily the root):
\begin{enumerate}
    \item Search the heap array for $k$. Since a heap has no directional search property, finding $k$ takes $O(n)$ time.
    \item Overwrite $k$ with the last element in the array and decrement the size.
    \item Repair the heap: if the replacement is greater than its parent, sift up; if smaller than either child, sift down. Repair takes $O(\log n)$.
    \item Total time: $O(n)$ due to linear search.
\end{enumerate}

\section{Complete C++ Implementation of a Max-Heap}

\begin{lstlisting}[style=cppstyle]
#include <iostream>
#include <vector>
#include <stdexcept>
#include <utility>
using namespace std;

class MaxHeap {
private:
    vector<int> heap;

    void siftUp(size_t i) {
        while (i > 0) {
            size_t parent = (i - 1) / 2;
            if (heap[parent] >= heap[i]) break;
            swap(heap[parent], heap[i]);
            i = parent;
        }
    }

    void siftDown(size_t i) {
        size_t n = heap.size();
        while (true) {
            size_t largest = i;
            size_t left = 2 * i + 1;
            size_t right = 2 * i + 2;

            if (left < n && heap[left] > heap[largest]) {
                largest = left;
            }
            if (right < n && heap[right] > heap[largest]) {
                largest = right;
            }
            if (largest == i) break;

            swap(heap[i], heap[largest]);
            i = largest;
        }
    }

public:
    // O(n) Bottom-Up constructor
    explicit MaxHeap(vector<int> values) : heap(std::move(values)) {
        for (size_t i = heap.size() / 2; i > 0; --i) {
            siftDown(i - 1);
        }
    }

    bool empty() const { return heap.empty(); }
    size_t size() const { return heap.size(); }

    // O(1) Max element inspection
    int getMax() const {
        if (heap.empty()) throw underflow_error("Heap is empty");
        return heap[0];
    }

    // O(log n) Insertion
    void insert(int value) {
        heap.push_back(value);
        siftUp(heap.size() - 1);
    }

    // O(log n) Extract Maximum
    int extractMax() {
        if (heap.empty()) throw underflow_error("Heap is empty");
        int maximum = heap[0];
        if (heap.size() == 1) {
            heap.pop_back();
            return maximum;
        }
        heap[0] = heap.back();
        heap.pop_back();
        siftDown(0);
        return maximum;
    }

    void display() const {
        for (int val : heap) cout << val << " ";
        cout << '\n';
    }
};

int main() {
    MaxHeap h(vector<int>{4, 10, 3, 5, 1});

    cout << "After construction: ";
    h.display();

    h.insert(15);
    cout << "After inserting 15: ";
    h.display();

    cout << "Maximum element: " << h.getMax() << '\n';
    cout << "Extracted maximum: " << h.extractMax() << '\n';

    cout << "After deleting maximum: ";
    h.display();
    cout << "Heap size: " << h.size() << '\n';

    return 0;
}
\end{lstlisting}

\subsubsection{Expected Output}
\begin{lstlisting}[style=pseudocode]
After construction: 10 5 3 4 1
After inserting 15: 15 5 10 4 1 3
Maximum element: 15
Extracted maximum: 15
After deleting maximum: 10 5 3 4 1
Heap size: 5
\end{lstlisting}

\section{Heap Complexity: Complete Comparison}

\begin{table}[htbp]
\centering
\small
\begin{tabularx}{\textwidth}{l X X}
\toprule
\textbf{Operation} & \textbf{Time Complexity} & \textbf{Algorithmic Reason} \\
\midrule
Find Maximum / Minimum    & $O(1)$                    & Root is always the extremum \\
Insert Element            & $O(\log n)$ worst-case    & Sifts up tree height $\lfloor \log_2 n \rfloor$ \\
Delete Max / Min (Root)   & $O(\log n)$ worst-case    & Sifts down tree height \\
Heapify Single Node       & $O(\log n)$ worst-case    & Floats down subtree height \\
Build Heap (Bottom-Up)    & $O(n)$                    & Leaf-heavy distribution of complete tree \\
Build Heap (Insertion)    & $O(n \log n)$             & Repeated insertions \\
Search for Arbitrary Key  & $O(n)$                    & No left/right directional ordering \\
Heap Sort                 & $O(n \log n)$             & $n$ extractions of $O(\log n)$ cost \\
Storage for $n$ elements  & $O(n)$                    & Compact 1D array \\
\bottomrule
\end{tabularx}
\caption{Comprehensive operational complexity reference for Binary Heaps.}
\end{table}

\subsection{Binary Heap versus Priority Queue}
A \textbf{priority queue} is an Abstract Data Type (ADT) defining operations (\texttt{push}, \texttt{top}, \texttt{pop}). A \textbf{binary heap} is the standard concrete data structure used to implement a priority queue efficiently.

In C++ STL:
\begin{lstlisting}[style=cppstyle]
#include <queue>
#include <vector>
#include <functional>

// Default: Max-Priority Queue
priority_queue<int> maxPQ;

// Min-Priority Queue
priority_queue<int, vector<int>, greater<int>> minPQ;
\end{lstlisting}

\section{Heap Sort}

Heap sort is an in-place, comparison-based sorting algorithm using a binary heap.

\subsection{Algorithm (Ascending Order)}
\begin{enumerate}
    \item Build a max-heap from the input array ($O(n)$).
    \item Swap $A[0]$ (the maximum) with the last element $A[n-1]$ in the active heap.
    \item Reduce active heap size by $1$.
    \item Sift down the new root ($A[0]$) across the active heap ($O(\log n)$).
    \item Repeat steps 2--4 until active heap size is $1$.
\end{enumerate}

\subsection{Properties of Heap Sort}
\begin{itemize}
    \item \textbf{Time Complexity:} $O(n \log n)$ in Best, Average, and Worst cases.
    \item \textbf{Auxiliary Space:} $O(1)$ in-place.
    \item \textbf{Stability:} \textbf{Not stable}. Swapping elements across the heap destroys the relative order of identical keys.
\end{itemize}

\section{Common Heap Errors and Examination Traps}

\begin{enumerate}
    \item \textbf{Assuming Heap Array is Sorted:} A heap is \emph{not} a sorted array. Only parent-child vertical ordering holds.
    \item \textbf{Confusing Heap with BST:} Left children are not necessarily smaller than right children in a heap.
    \item \textbf{Sifting in Wrong Direction:} Insertion requires \emph{sift-up}; root deletion requires \emph{sift-down}; bottom-up construction runs \emph{sift-down} starting from the last internal node.
    \item \textbf{Selecting Wrong Child in Sift-Down:} In a max-heap, you must swap with the \textbf{larger} of the two children. Swapping with the smaller child leaves the heap property violated.
    \item \textbf{Starting Bottom-Up at the Wrong Index:} Begin at $\lfloor n/2 \rfloor - 1$, not at $n-1$. Leaves have no children and are already valid heaps.
    \item \textbf{Confusing Memory Heap with Data Structure Heap:} The runtime memory heap used by \texttt{malloc} or \texttt{new} is an operating system memory management pool, entirely distinct from the binary heap data structure.
\end{enumerate}

\section{Practical Applications of Heaps}
\begin{itemize}
    \item \textbf{Operating System Schedulers:} Priority scheduling for CPU jobs.
    \item \textbf{Graph Algorithms:} Dijkstra's shortest path and Prim's Minimum Spanning Tree (using a min-priority queue).
    \item \textbf{Top-$K$ Selection Problems:} Maintaining the $k$ largest or smallest elements from a massive streaming dataset in $O(n \log k)$ time using a bounded heap.
    \item \textbf{Discrete Event Simulation:} Processing events chronologically based on timestamps.
    \item \textbf{Huffman Coding:} Constructing optimal prefix trees for data compression.
\end{itemize}

\section{Likely Recruitment Examination Questions}

\begin{table}[htbp]
\centering
\small
\begin{tabularx}{\textwidth}{p{7cm} X}
\toprule
\textbf{Question} & \textbf{Key Answer} \\
\midrule
What is a linked list? & A linear sequence of nodes connected via pointers. \\
Difference between singly and doubly linked list? & Singly has 1 pointer (\texttt{next}); doubly has 2 pointers (\texttt{prev}, \texttt{next}). \\
How does circular linked list differ from singly? & Last node points back to the head node instead of \texttt{nullptr}. \\
Time complexity of insertion at head? & $O(1)$ pointer update. \\
Why does linked list not support $O(1)$ random access? & Nodes are non-contiguous; reaching index $k$ requires sequential traversal. \\
What is a binary heap? & A complete binary tree satisfying a heap-ordering rule. \\
What is at the root of a max-heap? & The maximum element in the heap. \\
Formula for last non-leaf in 0-indexed heap? & $\lfloor n/2 \rfloor - 1$. \\
Formulas for left and right children (0-based)? & Left: $2i + 1$; Right: $2i + 2$. \\
Formula for parent index (0-based)? & $\lfloor (i - 1) / 2 \rfloor$. \\
Complexity of bottom-up heap construction? & $O(n)$ linear time. \\
Complexity of insertion into binary heap? & $O(\log n)$ worst-case. \\
Complexity of deleting the root from heap? & $O(\log n)$ worst-case. \\
Complexity of inspecting root maximum? & $O(1)$. \\
Can arbitrary key be searched in $O(\log n)$? & No; worst-case arbitrary search is $O(n)$. \\
Time complexity of Heap Sort? & $O(n \log n)$ in all cases (best, average, worst). \\
Is Heap Sort stable? & No, it does not preserve relative order of duplicates. \\
Is a max-heap a Binary Search Tree? & No; heap enforces vertical order, BST enforces horizontal order. \\
Why does bottom-up build take $O(n)$? & Most nodes reside near the leaves and travel very small distances. \\
\bottomrule
\end{tabularx}
\caption{Rapid-fire recruitment examination questions and answers.}
\end{table}

\section{Final Quick-Revision Sheet}

\subsection{Linked Lists Quick Revision}
\begin{itemize}
    \item \textbf{Singly Linked List:} Unidirectional (\texttt{next}); $O(1)$ insert at head; $O(n)$ arbitrary search/delete.
    \item \textbf{Doubly Linked List:} Bidirectional (\texttt{prev}, \texttt{next}); $O(1)$ delete of known node.
    \item \textbf{Circular Linked List:} Last node links to first; $O(1)$ front and end insertion with \texttt{tail} pointer.
    \item Traversal loop for circular list requires \texttt{do \{ ... \} while (curr != head);}.
\end{itemize}

\subsection{Max-Heaps Quick Revision}
\begin{itemize}
    \item Complete binary tree stored in array without pointers.
    \item Parent $\geq$ Children; Maximum at $A[0]$.
    \item 0-based formulas: $\text{left} = 2i + 1$, $\text{right} = 2i + 2$, $\text{parent} = \lfloor(i-1)/2\rfloor$.
    \item Build heap bottom-up: $O(n)$.
    \item Insert: append and sift up ($O(\log n)$).
    \item Delete Max: replace root with last element and sift down ($O(\log n)$).
    \item Get Max: $O(1)$.
    \item Heap Sort: $O(n \log n)$ time, $O(1)$ auxiliary space, unstable.
\end{itemize}

\begin{quote}
\textbf{Final Architectural Takeaway:} Linked lists optimize dynamic node insertion and deletion without shifting contiguous memory. Binary heaps leverage the structural completeness of binary trees to deliver $O(1)$ extremum access and $O(\log n)$ priority updates with zero pointer memory overhead.
\end{quote}
